\section{结论}\label{sec:concl}

本文围绕啤酒风味的三项任务，先解决"评分端点从哪来"，再分别给出预测、解释与决策三类答案。

\textbf{第一问}：以 2197 条逐人 0--6 评分为依据重建了 50 个感官目标的面板均分与七级分布。按酒三折五次验证下，固定等权融合在核心 30 项上的均分误差为 0.50222 分，低于训练均值 0.60450 分，但只比风格中位数基线（0.50969 分）低 0.00747 分，差值的 95\% 重采样区间为 $[-0.01914,0.00418]$，含 0。该比较未提供融合模型稳定优于已知风格中位数基线的充分证据。新增同参数ET对照中，已知风格加化学比仅风格ET的核心MAE低0.01426分，表明在该模型设置下化学输入仍有预测信息；相对风格中位数的差值区间跨0。七级分布模型的累计分布平方距离为 0.02178（总体基线 0.02820），面板众数一致率 80.52\%，但非零众数一致率只有 24.33\%：总体一致率主要由 78.03\% 的零等级贡献，精细等级识别仍是短板。保形区间逐属性平均覆盖 90.06\%、平均宽 1.7342 分，但按品评人数分组后，3--4 人组只有 74.00\%。缺失机制实验进一步显示，系统性低值缺失（MNAR）使核心误差增加 0.0273 分，约为随机缺失 20\%（0.0139 分）的两倍。

\textbf{第二问}：231 个字段在 $|r|\ge0.7$ 下构成 109 条边的共表达图，但 150 个字段是孤立节点。图正则模型在单折口径下略差于无图（0.537035 对 0.536911）、在聚合口径下略好（0.524865 对 0.526148，区间 $[-0.002341,-0.000229]$），最终全数据模型选得 $\lambda=0$，即 Top-3 实际来自普通岭回归。150 条候选中有 36 条通过"频率 $\ge0.5$、至少 5 折入选、置换贡献为正"的描述性条件；把 80 个化学异常单元格置为缺失后 Top-3 保留 97.64\%，仅使用38个逐点一致字段、61个单位登记字段和226个非综合字段的32项交集时降至36.18\%，该交集与32个异常字段不重合。36条为满足描述性支持条件的候选，另外114条为探索性候选；网络用于描述共变结构，候选不代表因果效应。

\textbf{第三问}：三种风格各生成 40 条谱。受约束局部凸组合在 Tripel 上的平均目标风格概率为 0.3358、目标预测比例 0.695，均高于重采样的 0.2992 与 0.565，但分布距离 MMD$^2$ 反而更大（0.1158 对 0.1034），说明"更像目标风格"与"更接近真实分布"不是一回事。迁移任务中，55 条"来源 $\times$ 划分"记录上搜索可行 51 条、独立代理支持 11 条、再加原始尺度 stale hops 约束后 6 条；真实留出的 10 条记录（7 款酒）依次为 10、2、1 条。严格协议（四项浓度逐项不增）下最大 Tripel 份额全为 0，且 55 条记录都找到了正余量不可行证书，阻塞主要来自 DMS 与反式-2-壬烯醛；参考阈值下30条记录仅满足化学正份额条件，联合风格与stale hops检查后没有候选通过J代理复核；2倍阈值敏感性下有9条重复划分记录获得J代理支持，其中3条为T内部留出来源，均未经过新批次制备和盲评。

\textbf{总体判断}：在现有数据上，风格基线具有较强预测能力；化合物候选需要结合共变模块和单项稳定性解释；缺陷阈值、代理模型分歧与真实制备共同决定迁移结果。后续工作应按此顺序推进：先补人员留出与真实品评，再谈更复杂的模型。
